\chapter{Scientific Computing: NumPy \& Pandas}

While standard Python lists are versatile, they are slow for massive numerical data because elements are pointers scattered in memory. The scientific Python ecosystem—led by \textbf{NumPy} and \textbf{Pandas}—provides contiguous memory layout, vectorization, and fast C-level execution for data manipulation.

\section{NumPy Arrays vs Python Lists}

\textbf{NumPy} (Numerical Python) introduces the \texttt{ndarray} (N-dimensional array), an array of homogeneous elements stored contiguously in memory.

\begin{definitionbox}{NumPy Arrays and Vectorization}
\textbf{Vectorization} allows mathematical operations to operate on entire arrays simultaneously without explicit Python \texttt{for} loops.
\begin{lstlisting}
import numpy as np

# Creating NDArrays
arr = np.array([1, 2, 3, 4, 5])

# Vectorized Arithmetic
print(arr * 2)       # Output: [ 2  4  6  8 10]
print(arr ** 2)      # Output: [ 1  4  9 16 25]

# Broadcasting: Matrix + Scalar
matrix = np.ones((2, 3))  # 2x3 matrix of ones
print(matrix + 5)
\end{lstlisting}
\end{definitionbox}

\begin{videocard}{IITM BS Lecture: Introduction to NumPy (Arrays vs Lists)}{https://www.youtube.com/watch?v=O0IrNk1qoUE}
Official lecture explaining memory layout, performance comparisons, and vectorization benefits.
\end{videocard}

\section{Pandas Series and DataFrames}

\textbf{Pandas} provides high-performance tabular data structures: the 1D \textbf{Series} (labeled column) and the 2D \textbf{DataFrame} (tabular dataset with labeled rows and columns).

\begin{theorembox}{Pandas DataFrame Mechanics}
\begin{lstlisting}
import pandas as pd

# Creating a DataFrame from a Dictionary
data = {
    "Name": ["Alice", "Bob", "Charlie"],
    "Age": [25, 30, 35],
    "Salary": [70000, 80000, 95000]
}
df = pd.DataFrame(data)

# Accessing Columns and Filtering Rows
print(df["Name"])
high_earners = df[df["Salary"] >= 80000]
print(high_earners)
\end{lstlisting}
\end{theorembox}

\textbf{Data Science Relevance:} Pandas is the default tool for Exploratory Data Analysis (EDA), CSV/Parquet loading (\texttt{pd.read\_csv()}), handling missing values (\texttt{df.fillna()}), and feature engineering.

\begin{videocard}{IITM BS Lecture: Why Pandas? Data Analysis Made Easy}{https://www.youtube.com/watch?v=HCxhHo8tV4U}
Official lecture introducing Pandas for data manipulation and tabular alignment.
\end{videocard}

\begin{videocard}{IITM BS Lecture: Pandas Series and DataFrames}{https://www.youtube.com/watch?v=aPjiMs7CR7s}
Official lecture detailing indexing, row selection (\texttt{.loc}, \texttt{.iloc}), and column operations.
\end{videocard}

\newpage
\section{Practice Exercises}
\textit{Detailed step-by-step solutions are available in the Appendix.}

\begin{enumerate}
    \item \textbf{NumPy Vectorization:} Given \texttt{arr = np.arange(1, 11)} (integers $1$ to $10$):
    \begin{enumerate}
        \item Compute a boolean mask for all odd numbers.
        \item Use the boolean mask to extract only the odd elements.
        \item Replace all even numbers in \texttt{arr} with $-1$.
    \end{enumerate}

    \item \textbf{Matrix Statistics with NumPy:} Create a $3 \times 3$ matrix with values $1$ through $9$ using \texttt{np.arange(1, 10).reshape(3, 3)}. Calculate:
    \begin{enumerate}
        \item Sum of all elements in the matrix.
        \item Column-wise means (\texttt{axis=0}).
        \item Row-wise maximums (\texttt{axis=1}).
    \end{enumerate}

    \item \textbf{Pandas Filtering \& Selection:} Given the DataFrame:
\begin{lstlisting}
import pandas as pd
df = pd.DataFrame({
    'City': ['Mumbai', 'Delhi', 'Bangalore', 'Chennai'],
    'Temp': [32, 40, 24, 34],
    'Humidity': [80, 45, 65, 75]
})
\end{lstlisting}
    Filter and print rows where \texttt{Temp > 30} AND \texttt{Humidity > 70}.

    \item \textbf{Pandas Missing Data Handling:} Create a DataFrame with missing values \texttt{np.nan}:
    \texttt{df = pd.DataFrame(\{'A': [1.0, 2.0, np.nan, 4.0], 'B': [np.nan, 10.0, 20.0, 30.0]\})}.
    Fill missing values in column \texttt{'A'} with its mean, and column \texttt{'B'} with \texttt{0.0}.

    \item \textbf{Data Science Application (Groupby Aggregation):} Given sales transaction data:
\begin{lstlisting}
sales_df = pd.DataFrame({
    'Category': ['Tech', 'Tech', 'Office', 'Office', 'Tech'],
    'Revenue': [1200, 800, 150, 300, 2000]
})
\end{lstlisting}
    Write a 1-line Pandas expression using \texttt{.groupby()} to calculate the total revenue per category.
\end{enumerate}